\chapter{Литературный аудит граничной трансгрессии класса \texorpdfstring{$1/7$}{1/7}}

Предыдущая глава получила точное равенство стабильных $K_0$-рангов, но
не оператор переноса. Настоящий гейт проверяет первичную литературу по
четырём близким механизмам: расширению Тёплица, mapping cone,
Kasparov-произведению и соединяющему гомоморфизму Милнора для хопфовых
расслоений.

Задача литературного аудита не состоит в поиске готовой физической теории.
Нужно установить, существует ли стандартный математический механизм,
который канонически переводит граничный класс в объёмный $K$-класс, и
какие дополнительные данные этот механизм требует.

\section{Расширение Тёплица}

Классическое расширение имеет вид
\begin{equation}
 0\longrightarrow\mathcal K
 \longrightarrow\mathcal T
 \longrightarrow C(S^1)
 \longrightarrow0.
 \label{eq:v5-one-seventh-toeplitz-extension}
\end{equation}
Его граничный морфизм
\begin{equation}
 \partial_{\mathcal T}:K_1(C(S^1))\longrightarrow K_0(\mathcal K)
\end{equation}
является индексом оператора Тёплица и переводит класс унитарной функции в
её число обёртываний с зависящим от ориентационной конвенции знаком.
Современный обзор этой конструкции и её роли в операторной $K$-теории
дан Аричи и Месландом, \path{arXiv:1911.05823}.

Тензорирование коэффициентной алгеброй сохраняет расширение:
\begin{equation}
 0\longrightarrow\mathcal K\otimes M_{105}
 \longrightarrow\mathcal T\otimes M_{105}
 \longrightarrow C(S^1)\otimes M_{105}
 \longrightarrow0.
 \label{eq:v5-one-seventh-coefficient-toeplitz}
\end{equation}
Пусть $u_H$ --- хопфова граничная унитарная функция с
$\operatorname{wind}(u_H)=1$, а
\begin{equation}
 q_0=p_{15}\otimes P_0,
 \qquad \rank q_0=15.
\end{equation}
Тогда естественный внешний продукт удовлетворяет
\begin{equation}
 \boxed{
 \partial_{\mathcal T}
 \bigl([u_H]\boxtimes[q_0]\bigr)
 =\pm[q_0]
 }
 \label{eq:v5-one-seventh-toeplitz-candidate}
\end{equation}
и имеет абсолютный ранг пятнадцать. Нормированный след равен
\begin{equation}
 \frac{15}{105}=\frac17.
\end{equation}
По предыдущему гейту этот класс совпадает по стабильному рангу с
\begin{equation}
 ([p_{20}]-[p_{15}])\otimes[I_3].
\end{equation}

\section{Почему граничный морфизм действительно каноничен}

В $KK$-теории граничная карта вычисляется Kasparov-произведением с
классом самого расширения. Бурн, Келлендонк и Ренни строят такой механизм
для bulk--edge соответствия и подчёркивают необходимость вещественной
версии при наличии антилинейных симметрий; см.
\path{arXiv:1604.02337}. Поэтому после фиксации расширения, ориентации
$S^1$, класса $u_H$ и коэффициентного проектора $q_0$ результат
\eqref{eq:v5-one-seventh-toeplitz-candidate} является каноническим на
уровне $K$-класса.

Но каноничность класса не означает каноничность отдельной матрицы или
частичной изометрии. $K$-теория факторизует по стабилизации и гомотопии.
Именно поэтому она способна устранить свободу $U(15)$ на уровне класса,
не выбирая единственный представитель в исходном конечном модуле.

\section{Mapping cone и условия применимости}

Кэри, Филлипс и Ренни используют mapping cone, чтобы связать нечётное
индексное спаривание с чётным спариванием при граничных условиях
Атьи--Патоди--Зингера; см. \path{arXiv:0711.3028}. Эта конструкция
подтверждает правильность общего языка перехода от границы к объёму, но
начинается с уже заданных включения алгебр и Kasparov-модуля. Она не
создаёт их из одного равенства следов.

Аричи и Ренни показывают, что для бимодулей конечного индекса точные
последовательности Cuntz--Pimsner и mapping cone могут быть явно
сопоставлены; классы mapping cone представляются стабильными классами
частичных изометрий, а недостающие карты строятся Kasparov-произведением;
см. \path{arXiv:1605.08593}. Однако их механизм требует
$C^*$-соответствия над одной коэффициентной алгеброй или предварительного
расширения скаляров. Наш $M_{20\times15}$ является соответствием между
двумя различными углами, но его самосоответствие, алгебра Фока и
Cuntz--Pimsner расширение ещё не заданы.

\section{Свежий хопфово-милноровский результат}

Работа Д'Андреа, Хаяца, Мащика и Зелиньского, версия 3 от 29 апреля
2026 года, строит явный соединяющий гомоморфизм Милнора в контексте
Hopf--Galois расслоений и выражает clutching-класс через матрицу
представления; см. \path{arXiv:2512.08304v3}. Она особенно близка нашей
постановке по трём причинам:
\begin{enumerate}
 \item исходным объектом служит хопфово расслоение;
 \item граничные данные порождают явный идемпотент;
 \item результат лежит в положительном конусе $K_0$ в исследованных
 квантовых примерах.
\end{enumerate}

Но последний пункт авторы прямо характеризуют как квантовый эффект,
отсутствующий в классическом случае. Поэтому их положительный
представитель нельзя автоматически перенести в текущий классический
$M_{35}$. Для применения формулы нужно сначала построить конкретную
pullback/Hopf--Galois алгебру проекта и её clutching-матрицу.

\section{Что литература закрывает}

Литература позволяет строго утверждать следующее.
\begin{enumerate}
 \item Граничный класс обёртывания действительно может канонически
 перейти в чётный $K_0$-класс через расширение Тёплица.
 \item Коэффициентный проектор переносится внешним произведением, поэтому
 ранг пятнадцать и вес $1/7$ сохраняются без выбора базиса.
 \item Граничная карта является произведением с классом расширения, а не
 следствием одного равенства рангов.
 \item Mapping cone и Cuntz--Pimsner дают подходящий язык для
 операторного представителя, но требуют заранее заданного соответствия и
 расширения.
 \item Хопфово-милноровская конструкция может дать явный идемпотент, но
 требует проектно-специфической pullback-алгебры.
\end{enumerate}

\section{Что в проекте ещё отсутствует}

Текущий родитель $M_{35}$ не содержит алгебру Тёплица, пространство Харди
или выбранное заполнение окружности диском. Хопфова линия и ориентация
$L/L^*$ дают класс $u_H$ и знак обёртывания, но ещё не доказывают, что
физический родитель обязан реализовывать именно расширение
\eqref{eq:v5-one-seventh-coefficient-toeplitz}.

Также не построены:
\begin{enumerate}
 \item самосоответствие Cuntz--Pimsner над одной алгеброй;
 \item проектная pullback-алгебра для формулы Милнора;
 \item вещественный $KKO$-подъём, совместимый с KO6;
 \item оператор Дирака, энергия и локализованный индекс на физическом
 носителе.
\end{enumerate}

\begin{nogocriterion}[Запрет подмены класса физическим оператором]
Формула Тёплица выводит канонический $K_0$-класс только после задания
класса расширения. Она не доказывает, что конечный родитель $M_{35}$ сам
порождает пространство Харди, динамический оператор или положительный
ранг-пять проектор в $H_{20}$.
\end{nogocriterion}

\section{Вердикт и следующий тест}

Общий механизм полностью известен и не является новизной проекта.
Положительный результат состоит в том, что найденная нами арифметика
$5\cdot3=15\cdot1$ точно совместима с коэффициентным индексным
морфизмом Тёплица. Проектная новизна могла бы состоять только в выводе
класса расширения из уже существующих хопфовой линии и соответствия
Мориты.

Следующий гейт должен явно построить
\eqref{eq:v5-one-seventh-coefficient-toeplitz}, вычислить его граничную
карту и проверить KO6. Если алгебра Тёплица или её поляризация остаётся
внешним выбором, ветвь сохраняется как точная топологическая модель, но не
как вывод родительской физики.

\begin{protocol}
\begin{enumerate}
 \item общий механизм граничной трансгрессии: \textbf{известен};
 \item естественная формула Тёплица для текущих классов:
 \textbf{получена};
 \item абсолютный ранг образа: \textbf{пятнадцать};
 \item нормированный вес образа: \textbf{$1/7$};
 \item устранение свободы представителя на уровне $K_0$:
 \textbf{получено};
 \item расширение Тёплица внутри родителя $M_{35}$:
 \textbf{не выведено};
 \item проектная формула Милнора: \textbf{требует новой pullback-алгебры};
 \item KO6/KKO-подъём: \textbf{не выполнен};
 \item физический оператор и энергия: \textbf{не получены};
 \item следующий гейт: \textbf{явная коэффициентная карта Тёплица}.
\end{enumerate}
\end{protocol}

Машинный сертификат сохранён в
\path{s2t_v5_one_seventh_boundary_transgression_literature_gate_results.json}.